\documentclass[../../main2.tex]{subfiles}

\begin{document}

	% ======================================================================
	% 骨架文件：小节标题与追问链为终版；正文由后续会话填充。
	% 扩增模式标注：本章 = 有压扩增（稀缺约束）。
	% 原书材料接入点：
	%   - 原书 part3 ch1（经济基础）→ 本章主体，改由「文化+稀缺」两步推导导出
	%   - §10.1 是本 Part 的枢纽节（notes/00-architecture.md 明示「必须完整写」）
	% ======================================================================

	\chapter{通道二：经济（模因 + 稀缺约束）}
	\label{ch:economy}

	\begin{motivation}
		上一章遗留的问题：模因传播不消耗原体、也就无价——但现实里大量传播必须付费。如果一种模因的传播\textbf{必须}消耗稀缺资源，它还是模因吗？\\
		\textbf{核心动机}：用「加入稀缺约束」这一步，从文化模因中推导出经济——供求不是经济学的基本假设，而是「模因 + 稀缺」两步推导的产物。\\
		\textbf{扩增模式}：有压——复制变转移，总量守恒，零和开始起作用。
	\end{motivation}

	\begin{logicmap}
	\begin{tikzpicture}[node distance=0.9cm, every node/.style={draw, rectangle, rounded corners, align=center}]
	\node (q) {传播要花钱，还是模因吗？};
	\node (s) [below=of q] {加入稀缺约束（§10.1，枢纽节）};
	\node (sd) [below=of s] {稀缺→供方；偏好→需求（§10.2）};
	\node (h) [below=of sd] {追涨杀跌的理性化（§10.3）};
	\node (c) [below=of h] {经济的贡献分解（§10.4）};
	\node (n) [below=of c] {交换真的等价吗？（抛给第11章）};
	\draw[-{Stealth}] (q) -- (s);
	\draw[-{Stealth}] (s) -- (sd);
	\draw[-{Stealth}] (sd) -- (h);
	\draw[-{Stealth}] (h) -- (c);
	\draw[-{Stealth}] (c) -- (n);
	\end{tikzpicture}
	\end{logicmap}

	% ==================================================================
	\section{加入稀缺约束}
	\label{sec:econ-scarcity}

	如果传播不消耗原体——你把一个想法告诉我，你手里还有那个想法——扩张就没有上限。一个笑话讲给一百个人听，它还是一百个笑话，不多不少。这是上一章说的无压扩增：复制 $A \to A+B$，原体保留，总量只增不减。

	但现实里大量传播必须付费。你买一杯咖啡，钱离开你的口袋进入商家的账户；你买一股股票，那笔资金从买方转移到卖方。这一步把复制变成了转移：$A \to B$，原体不再保留，总量锁死。守恒律登场——在这个通道里，价值不凭空创造也不凭空消灭，只是换了持有者。

	守恒带来了什么新性质？最直接的：零和开始起作用。在模因通道里，你多知道一个事实，我不因此少知道——双赢没有代价。但在经济通道里，你多赚一块钱，意味着某个人少赚了一块（或者多付了一块）。不是每一笔交易都零和——贸易可以做大蛋糕——但守恒律保证了蛋糕的总质量有账可查，而零和的部分（存量博弈、投机、租金提取）始终存在，并且往往是最激烈的部分。

	这就是经济行为从文化行为中分离出来的时刻。同一个模因——「这只股票会涨」——在朋友圈免费扩散是无压的；但它一旦驱动真金白银的买卖，就进入了有压通道，受守恒律约束。

	\begin{readingnote}
		你有没有想过，为什么「分享」这个词在朋友圈免费，在股市要收印花税？因为朋友圈走的是复制通道（我转发一条消息，原文还在），股市走的是转移通道（我卖出一股，那笔钱就不在我这里了）。两种通道的数学结构不同，所以规则不同。
	\end{readingnote}

	\paragraph*{逻辑衔接}
	守恒性落地之后，最古老的经济学图景自己长出来了。下一节做两步推导：稀缺产生供方，偏好产生需求。

	% ==================================================================
	\section{稀缺到供方，偏好到需求}
	\label{sec:econ-supplydemand}

	上一节建立了守恒律。现在做两步推导，把供求曲线从假设变成结论。

	第一步：稀缺产生供方。什么东西稀缺，什么就有人愿意提供——因为稀缺意味着持有者有议价权。供方不是「生产」的别名，而是「持有稀缺资源并等待交换」的行为集合。你囤了一箱口罩，你就是供方；你写了一个开源库，你不是供方（因为复制不消耗原体，无稀缺）。

	第二步：偏好产生需求。需求来自第2章建立的目标函数——个体在多个目标之间权衡后的选择倾向。你「需要」咖啡不是因为咖啡因（那是生理约束），而是因为「清醒地工作」这个目标在当前权重最高。偏好是模因层面的东西：它被文化通道塑造，但通过经济通道表达。

	价格从哪来？价格是压缩进一个数字的供需信息。买方和卖方各自携带大量局部信息（库存、产能、偏好强度、替代选项），这些信息的完整版本不可能全部传递——就像变长编码一样，价格是最短的那个码字，把「当前有多稀缺」和「当前有多想要」压成一个数。

	但这个压缩会失真。外部性是最典型的失真形式：工厂排污，污染成本不在价格里——价格只反映了买卖双方的私人信息，没反映对第三方的影响。市场失灵不是市场的bug，而是信息压缩的固有代价：任何压缩都有损，价格也不例外。

	\begin{questionbox}
		「供求定律有反例——吉芬商品越贵越买，炫耀性消费越贵越想要。推导还有效吗？」有效。反例不推翻「两步推导」，只说明偏好维度比单维更复杂。吉芬商品的「反常」来自偏好结构（面子效用 > 价格效用），恰好回到第2章的多目标框架——目标函数多了一个维度，供求曲线的形状就变了，但「稀缺→供方、偏好→需求」的两步推导本身没有失效。
	\end{questionbox}

	\paragraph*{逻辑衔接}
	价格是信号，但信号有噪声。噪声下的行为是理性的还是疯狂的？下一节「追涨杀跌的理性化」回答这个问题。

	% ==================================================================
	\section{追涨杀跌的理性化}
	\label{sec:econ-herding}

	追涨杀跌看起来非理性，但它有一个精确的理性内核。

	核心机制叫信息级联（information cascade）：当你不确定某件事的真假时，观察前面的人做了什么，比自己做独立判断更便宜——尤其是当你的私有信息很弱的时候。第一个人凭自己的信息买入；第二个人看到第一个人买了，结合自己微弱的信号，也买；第三个人看到前两个人都买了，即使自己觉得「好像不太对」，也会跟着买——因为前两个人的联合信息权重超过了自己的直觉。级联一旦形成，后来者的行为就不再携带新信息，全部变成对前人的模仿。

	泡沫是级联的上行端：价格持续上涨本身成为「有人愿意出更高价」的信号，吸引更多人加入，价格进一步上涨。崩盘是级联的下行端：一个足够大的抛售打破了「还会涨」的预期，级联反向运行，恐慌性抛售让价格雪崩。两端都是同一个机制——模仿作为信号处理——只是方向相反。

	这里有一个反直觉的结论：参与级联的每一个个体都可以是理性的。每个人都在做信息约束下的最优近似——「别人都买了我不买，要么我错了，要么他们知道我不知道的事」。但理性个体的加总产生了集体非理性：所有人都做局部最优，整体却偏离了基本面。这不是「群体愚蠢」，而是信息结构的固有局限——级联在信息不对称严重时是最优策略，但它的宏观后果是泡沫和崩盘。

	\begin{readingnote}
		下次你看到「全网都在买这个」的时候，可以问自己一个问题：我是有独立信息支撑我的判断，还是我只是在模仿前面的信号？如果是后者，你正在参与一个级联——这不一定错，但你应该知道自己处于级联里。
	\end{readingnote}

	\begin{reader_anticipation}
		\item 级联什么时候断裂？→ 第11章 §11.2：认同的凭空增减——政治守恒性弱于经济的根源。
	\end{reader_anticipation}

	\paragraph*{逻辑衔接}
	经济通道的机制齐了。它在结果里占多少？下一节做贡献分解——并直视一个常见的夸大。

	% ==================================================================
	\section{经济的贡献分解}
	\label{sec:econ-contribution}

	把经济通道的贡献写进分解公式：$C_{\\text{经济}}$ 衡量的是「因为稀缺约束和守恒律而改变的行为」在总行为中的份额。

	但这里必须诚实：经济贡献经常被高估，因为价值是测量的，不是物理的。价格会失真（上一节的外部性），GDP 不是自然常数——它依赖你选择测量什么、怎么加总。用价格做代理指标时，必须并列说明失真的方向和幅度，否则数字本身就是误导。

	被高估的典型情形：金融交易金额巨大，但很多只是存量资产的换手，不创造新价值。被低估的典型情形：家务劳动、生态服务、无偿照护——这些有真实经济价值但不进入价格体系，所以不出现在 GDP 里。\lowfreq{测量边界}：任何基于价格的贡献分解，其误差范围取决于未定价行为的占比。

	\paragraph*{逻辑衔接}
	经济交换在课本里是等价交换——双方地位对称，一手交钱一手交货。真的吗？一旦交换双方的身份变得重要——你是老板我是员工、你是买方我是卖方、你有关系我没有——等价交换就变形了。下一章从这里开始。

	% ==================================================================
	\section{本章小结与衔接}
	\label{sec:econ-summary}

	经济 = 加了稀缺约束的模因行为。一个想法在朋友圈扩散是无压的复制；同一套逻辑驱动真金白银的买卖时，就变成了有压的转移，受守恒律约束。守恒带来了零和博弈的可能，稀缺产生了供方，偏好产生了需求，价格是供需信息的压缩编码，级联是不完全信息下模仿作为信号的理性策略。文化 $\\supset$ 经济：经济通道是文化通道的一个子集，加了守恒约束的那个。

	\paragraph*{逻辑衔接}
	第11章「通道三：政治（经济 + 身份权重）」：当交换中加入不可货币化的偏好——人情、义务、权力——政治从经济中分离出来。

\end{document}
